\section{Case Study Design}
\label{sec:case-study}

The case-study design follows the logic of the prior PCI-DSS inspired experiment:
compare environmental assessment under flat and segmented network conditions. The
new study adds machine-readable evidence artifacts so that the same type of
scenario can be assessed by deterministic scripts, language-model agents, and a
watcher-enabled multi-agent workflow.

Each case is represented by:
\begin{itemize}[leftmargin=*]
  \item an asset inventory;
  \item vulnerability findings;
  \item network topology;
  \item firewall rules;
  \item business-impact mappings;
  \item PCI scope declarations; and
  \item expected expert labels used only for evaluation.
\end{itemize}

The experimental comparison includes four conditions:
\begin{enumerate}[leftmargin=*]
  \item historical human-baseline findings from the prior study;
  \item a deterministic rule-based baseline;
  \item a language-model baseline without local tool access; and
  \item the AI Bridge watcher workflow with local evidence access.
\end{enumerate}

The main outcomes are exact match against expert labels, environmental score
error, correct direction of before-after change, reproducibility across repeated
runs, and evidence coverage. Evidence coverage is defined as the proportion of
scoring decisions that can be traced to a concrete artifact such as an asset
record, topology entry, firewall rule, PCI scope declaration, or business-impact
statement.
